%!TeX program=xelatex
\documentclass[11pt,a4paper]{tutorial}

\usepackage{enumitem,booktabs,array}
\usepackage{multicol}

\newcommand{\Tr}{\operatorname{Tr}}
\newcommand{\Lang}{\mathcal{L}}
\newcommand{\EX}{\mathrm{EX}}
\newcommand{\AX}{\mathrm{AX}}
\newcommand{\EF}{\mathrm{EF}}
\newcommand{\AF}{\mathrm{AF}}
\newcommand{\EG}{\mathrm{EG}}
\newcommand{\AG}{\mathrm{AG}}

\setlist[enumerate]{label=(\alph*),itemsep=3pt,topsep=3pt}
\setlength{\parindent}{0pt}

\semester{Z. 2026}
\coursetitle{Concurrency Theory}
\lecturer{P. Hofman}

\tutorial{16}{10}{2026}

\documenttitle{Tutorial 3, \formatdate, by \textsc{vrm}}
\abstract{Languages, simulation, and logical abstraction}

\renewcommand{\thesection}{\Alph{section}.}

\begin{document}
\maketitle 

\paragraph{Conventions.}
For an action-labelled LTS, let
\[
 \Lang(s)=\{w\in\Sigma^*:s\xrightarrow{w}s'\text{ for some }s'\}
\]
be its language of all finite traces, including $\varepsilon$.
For a Kripke structure, $\Tr(s)$ denotes its set of infinite words of
state labels; the transition relation is total. LTL uses universal path
semantics. Write $s\preceq t$ when $t$ simulates $s$; on Kripke
structures, simulations preserve state labels exactly.
HML uses action modalities on LTSs, while CTL uses state propositions
on Kripke structures.\par

\paragraph{Scope.} Give explicit witnesses and short arguments. Do not compute
maximal bisimulation relations or use partition refinement.
Exercises marked $\star$ are optional extensions.


\section{What does an abstraction preserve?}
\begin{exercise}[Equal languages, different branching]
Construct two finite pointed LTSs $(T,s)$ and $(U,t)$ such that
$\Lang(s)=\Lang(t)$ but $s\not\sim t$.
\begin{enumerate}
\item Give an HML formula true at one initial state and false at the other.
\end{enumerate}
\end{exercise}

\begin{exercise}[Bisimilar but not isomorphic]
Construct two finite pointed LTSs with different numbers of reachable
states whose initial states are bisimilar. Exhibit a bisimulation witnessing
this and explain why the systems are not isomorphic.
\end{exercise}

\begin{exercise}[Simulation versus language inclusion]
\begin{enumerate}
\item Prove that $s\preceq t$ implies $\Lang(s)\subseteq\Lang(t)$.
\item Construct finite LTSs showing that the converse fails. Explain the
failure by identifying a transition that cannot be matched by a simulation.
\item Prove that, for deterministic LTSs, the converse holds.
Here deterministic means that each state has at most one successor
for each action.
\end{enumerate}
\end{exercise}

\begin{exercise}[CTL fragments and preservation under simulation]
Assume $s\preceq t$ in a Kripke structure. For each formula below, determine
whether the following implications hold for all such pairs:
\[
 (\rightarrow)\quad s\models\varphi\Rightarrow t\models\varphi,
 \qquad
 (\leftarrow)\quad t\models\varphi\Rightarrow s\models\varphi.
\]
Answer \emph{both}, \emph{only $\rightarrow$}, \emph{only $\leftarrow$},
or \emph{neither}. 
%\[
%\begin{array}{lll}
%\text{(i)}\ p\land\neg q & \text{(ii)}\ \EX p & \text{(iii)}\ \AX p\\[3pt]
%\text{(iv)}\ \EF(p\land\EX q) & \text{(v)}\ \AG(p\lor\AF q)
% & \text{(vi)}\ \AG\EF p\\[3pt]
%\text{(vii)}\ \EX p\land\AX q & \text{(viii)}\ \neg\EX\neg p
% & \text{(ix)}\ \EG p\lor\AF q.
%\end{array}
%\]
\begin{multicols}{3}
\begin{enumerate}[label=(\roman*)]
  \item \(p\land\neg q\)
  \item \(\EX p\)
  \item \(\AX p\)
  \item \(\EF(p\land\EX q)\)
  \item \(\AG(p\lor\AF q)\)
  \item \(\AG\EF p\)
  \item \(\EX p\land\AX q\)
  \item \(\neg\EX\neg p\)
  \item \(\EG p\lor\AF q\)
\end{enumerate}
\end{multicols}
Justify valid directions using preservation results. For each invalid
direction, give a finite counterexample and a witnessing simulation.
\end{exercise}

\begin{exercise}[Trace equivalence and a distinguishing CTL formula]
Consider the following two Kripke structures. Unlisted transitions are absent.
\begin{itemize}[itemsep=3pt,topsep=3pt]
\item $K$: $s\to x$, $x\to u$, $x\to v$, $u\to u$, $v\to v$;\\
$L(s)=L(x)=\varnothing$, $L(u)=\{p\}$, $L(v)=\{q\}$.
\item $H$: $t\to y$, $t\to z$, $y\to u'$, $z\to v'$, $u'\to u'$, $v'\to v'$;\\
$L(t)=L(y)=L(z)=\varnothing$, $L(u')=\{p\}$, $L(v')=\{q\}$.
\end{itemize}
\begin{enumerate}
\item Show that $\Tr(s)=\Tr(t)$ and conclude that $s$ and $t$ satisfy
the same LTL formulas.
\item Find a CTL formula distinguishing $s$ from $t$. Your formula should
express that one immediate successor offers two different possible futures.
\end{enumerate}
\end{exercise}

\begin{exercise}[Choose a sufficient abstraction]
For each verification task, choose a sufficient relation between a concrete
state $c$ and an abstract state $a$. State the direction of any trace
inclusion or simulation, and justify your answer in one sentence.
\begin{enumerate}
\item Transfer the LTL formula $\mathrm{F}p$ from $a$ to $c$.
\item Transfer the CTL formula $\AF p$ from $a$ to $c$. How does this task
relate to (a)?
\item Transfer the CTL formula $\EF p$ from $c$ to $a$.
\item Preserve the truth of the CTL formula $\EF\AG p$ in both directions.
\end{enumerate}
Use trace inclusion/equality, simulation, or bisimulation as appropriate.
In (a)--(c), consider whether trace inclusion alone suffices; your relation
need only be sufficient, so do not claim necessity without proof.
\end{exercise}

\section{Language problems for finite and infinite-state systems}
\paragraph{Acceptance conventions.}
Write $L(A)$ for an \emph{accepted} finite-word language, distinguished from
the all-traces language $\Lang$ above. NFAs and PDAs accept by final state
after consuming the whole input; PDAs may have $\varepsilon$-moves.
For the process and counter models below, every transition consumes one
letter; there are no silent transitions.
\begin{itemize}[itemsep=3pt,topsep=3pt]
\item A \textbf{BPA} has finitely many rules $X\xrightarrow{a}\alpha$,
where $\alpha$ is a word of variables. Its configurations are words,
with $X\beta\xrightarrow{a}\alpha\beta$. An initial word is specified;
acceptance is by reaching the empty word.
\item A \textbf{BPP} has the same rule form but configurations are
multisets of variables. A rule replaces one occurrence of $X$ by
the multiset $\alpha$. An initial multiset is specified; acceptance
is by reaching the empty multiset.
\item A \textbf{VASS} has configurations $(q,\mathbf v)$ with
$\mathbf v\in\mathbb N^d$ and transitions
$q\xrightarrow{a,\mathbf z}q'$, enabled when
$\mathbf v+\mathbf z\geq\mathbf0$. An initial configuration and final
control states are specified; final counter values are unrestricted.
A deterministic VASS has at most one outgoing transition per
control state and letter.
\end{itemize}

\begin{exercise}[NFA inclusion and equivalence]
\begin{enumerate}
\item Give algorithms deciding $L(A)\subseteq L(B)$ and $L(A)=L(B)$
for two NFAs. Explain why simply interchanging final and non-final
states of $B$ does not compute its complement.
\item Show how to extract a distinguishing word when equivalence fails.
\item Give an upper bound on the length of a shortest distinguishing word
in terms of the numbers of states of $A$ and $B$.
\end{enumerate}
\end{exercise}

\begin{exercise}[PDA languages and regular specifications]
Let $P$ be a PDA and $A$ an NFA over the same alphabet.
\begin{enumerate}
\item Show that $L(P)\subseteq L(A)$ is decidable using closure
properties and an emptiness test.

\item Investigate whether $L(A)\subseteq L(P)$ is decidable.
Consider the special case $L(A)=\Sigma^*$.
(Show that halting problem for automaton with two pushdowns is undecidbale. Try to accept all incorrect runs of an automaton with two pushdowns.)
\item Determine whether inclusion and equivalence of two PDA languages
are decidable. Give reductions from a standard undecidable problem;
do not infer inclusion from equivalence or conversely without an argument.
\end{enumerate}
\end{exercise}

\begin{exercise}[A unary alphabet]
\begin{enumerate}
\item Prove that every unary context-free language is regular.
You may use the effective version of Parikh's theorem.
Deduce decidability of inclusion and equivalence for unary PDA languages.
\item Show that the all-traces language of \emph{any} unary LTS is either
$\{\varepsilon,a,\ldots,a^n\}$ for some $n\geq0$, or $a^*$.
\item For unary VASSs, give an effective procedure deciding inclusion and
equivalence of all-traces languages. You may use decidability of VASS
termination, meaning absence of an infinite run from the initial configuration.
\item Explain why (b) does not establish regularity of every unary
\emph{accepted {\bfseries language}} of an arbitrary infinite-state system.
\end{enumerate}
\end{exercise}

\begin{exercise}[Regularity and unbounded configurations]
Consider a one-counter VASS with initial configuration $(r,0)$,
final control state $f$, and transitions
\[
 r\xrightarrow{a,+1}r,\qquad
 r\xrightarrow{b,-1}f,\qquad
 f\xrightarrow{b,-1}f.
\]
\begin{enumerate}
\item Describe its accepted language and prove that it is not regular.
\item Describe its all-traces language and determine whether it is regular.
\item Construct a VASS with infinitely many reachable configurations
whose accepted language is regular. Explain why unboundedness alone
does not decide language regularity.
\end{enumerate}
\end{exercise}

\begin{exercise}[Deterministic VASSs and regular specifications]
Let $V$ be a VASS and $A$ an NFA.
\begin{enumerate}
\item Show that $L(V)\subseteq L(A)$ is decidable by a product construction.
State whether the required VASS test is reachability or coverability.
\item Now let $V_1,V_2$ be deterministic VASSs.
Design a procedure for deciding $L(V_1)\subseteq L(V_2)$.
Separate rejection by $V_2$ caused by a missing transition, an insufficient
counter value, or a non-final control state.
You may use decidability of VASS reachability to semilinear target sets.
\item Deduce decidability of accepted-language equivalence for this
deterministic model. 
\end{enumerate}
\end{exercise}

\begin{exercise}[$\star$ A decidability map: BPA, BPP, and VASS]
Complete the following table for the \emph{accepted-language conventions}
fixed above. Write D (decidable) or U (undecidable).
For each non-elementary result, supply a precise theorem reference and
check its acceptance convention. A full proof is not required.
\begin{center}
\begin{tabular}{lc|c|c}
\toprule
Model & $L(A)\subseteq L(B)$ & $L(A)=L(B)$ & Is $L(A)$ regular?\\
\midrule
NFA & & &\\
\hline
PDA & & &\\
\hline
BPA & & &\\
\hline
BPP & & &\\
\hline
VASS & & &\\
\hline
Deterministic VASS & & &\\
\bottomrule
\end{tabular}
\end{center}
\begin{enumerate}
\item Investigate which table entries change when accepted languages are
replaced by all finite trace languages. State each convention explicitly.
\end{enumerate}
\end{exercise}

\end{document}
